\section{Introduction}
\label{sec:introduction}

We study finite-particle approximation for nonconvex optimization over probability measures,
\begin{equation}
    \min_{\mu\in\operatorname{dom}F\subseteq\Ptwo(\R^d)} F(\mu),
    \label{eq:measure-optimization}
\end{equation}
where \(\Ptwo(\R^d)\) denotes the set of Borel probability measures with finite second moment;
the domain \(\operatorname{dom}F\) for our moment models is specified in \Cref{def:moment-objective}.
The canonical first-order dynamics is the Wasserstein gradient flow
\begin{equation}
    \partial_t \mu_t + \nabla\cdot(v_t\mu_t)=0,
    \qquad
    v_t(x)=-\dW F(\mu_t,x),
    \label{eq:intro-wgf}
\end{equation}
which is gradient descent in the \(2\)-Wasserstein geometry.
In mean-field models, \eqref{eq:intro-wgf} is the population equation used to
describe interacting-particle dynamics in the infinite-particle limit.
This viewpoint underlies generative modeling~\citep{arbel2019,chu2019pfd},
sampling~\citep{jordan1998}, and mean-field neural-network
training~\citep{chizat2018,songmei2018,nitanda2022convex}.
Related particle variational methods include Stein variational
gradient descent~\citep{liu2016}.
In several convex, displacement-convex, and entropy-regularized settings,
WGF and its particle approximations admit strong convergence and propagation-of-chaos
\citep{sznitman1991} guarantees~\citep{chen2025uniform,
suzuki2023uniformintime,nitanda2025propagationensemble}.
Global convergence guarantees are also available for particle dual averaging
in mean-field neural-network optimization~\citep{nitanda2022pda}.

For nonconvex objectives, however, first-order stationarity does not rule out
saddle points. Such second-order landscape structure already appears in
mean-field models, including the attention model studied by
\citet{kim2024nonlinear}.
Inspired by perturbed gradient descent~\citep{GeHuangJinYuan2015,pmlr-v75-jin18a},
\citet{yamamoto2025hessianguided} introduce Hessian-guided perturbed Wasserstein
gradient flow (PWGF), which perturbs the population measure along
Hessian-guided Gaussian vector fields to escape saddles. Their guarantees
concern the population measure flow. A finite implementation replaces both the
drift and the Hessian-guided jump covariance by empirical quantities,
introducing errors along the flow and at each perturbation.
The interacting particles are generally dependent, so standard sampling bounds
for independent samples do not apply directly.
\Needspace{4\baselineskip}
We therefore ask:
\begin{center}
\emph{When can a finite system of interacting particles accurately approximate population PWGF?}
\end{center}

We answer this finite-sample question through the following contributions.

\Needspace{5\baselineskip}
\maintextheading{Contributions.}
\begin{enumerate}[leftmargin=1.35em,label=\arabic*.,itemsep=0.01em,topsep=0.05pt,parsep=5pt]
    \item \textbf{Tracking and objective values.}
    When the reference path satisfies a controlled-curvature path condition,
    we obtain particle and objective-value tracking bounds throughout \([0,T]\)
    under regularity assumptions and a prescribed deterministic common jump schedule
    (\Cref{theorem-main-r_t}).
    The bounds retain their \(T\)-dependence and weight ODE residuals and jump
    mismatches by separate curvature-response factors. By retaining the
    accumulated curvature, they capture how later positive curvature damps
    earlier errors. For fixed horizons, the ODE and jump prefactors are
    \(N^{-1/2}\) and \(N^{-1/4}\), respectively, up to arbitrarily small powers of \(N\).
    A variance-plus-cosine example gives a growing-horizon tracking guarantee
    (\Cref{ex:variance-cosine-curvature-recovery}).

    \item \textbf{A population-first coupling for Gaussian jumps.}
    At each jump, we construct a finite-dimensional standard Gaussian vector that
    reproduces the population-field evaluations and drives the empirical update.
    This preserves the correct marginal dynamics and the reference particles'
    conditional i.i.d. law given the population fields, while reducing the jump
    comparison to a covariance-factor mismatch
    (\Cref{subsec:auxiliary-particles,subsec:proof-sketch-after-assumptions}).

    \item \textbf{Local geometry of matrix factorization.}
    Matrix factorization provides a product model for the nonlinear-moment
    geometry arising in matrix and tensor decompositions. In a regularized
    product-measure model, we establish attraction along an anti-symmetric
    branch, descent transverse to that branch, and positive second variation
    in two regions under spherical product pushforward perturbations
    (\Cref{sec:matrix-factorization}). These local results motivate the
    positive--negative--positive curvature scenario.
\end{enumerate}

\maintextheading{Relation to prior work.}
\citet{yamamoto2025hessianguided} analyze population PWGF escape and convergence,
but do not bound its population-to-empirical approximation error.
\citet{glasgow2025meanfield} track continuous projected neural-network dynamics
using local particle curvature, without Hessian-induced Gaussian jumps.
Reference-particle couplings and Gaussian covariance-factor alignment have
precedents~\citep{sznitman1991,masarotto2019procrustes}. Here the coupling must
preserve the prescribed population fields, because the conditional
reference-particle law given those fields is the input to concentration.
We establish this compatibility and control the resulting jump covariance
residual through weighted features, then propagate continuous and jump errors
through the sampled curvature.

\maintextheading{Scope.}
The population PWGF proposed by \citet{yamamoto2025hessianguided}
selects perturbation times adaptively.
Prescribed schedules provide a reference point for studying how adaptive
timing affects finite-particle approximation.
Our results quantify finite-particle approximation errors under a common
prescribed schedule while retaining state-dependent drift and jump covariance.
The tracking bound applies along reference paths satisfying
the controlled-curvature path condition in \Cref{theorem-main-r_t},
which bounds error propagation along those paths.
The variance-plus-cosine example verifies this condition with high probability
and gives a growing-horizon tracking guarantee with a single prescribed
jump (\Cref{ex:variance-cosine-curvature-recovery}).
Extending this verification to broader models and synchronizing adaptive
jump times remain open directions.
